\section{Conclusion}
\label{sec:conclusion}

This review has mapped the methodological landscape and the
application landscape of surrogate models in optimisation-based
energy-system modelling, with an explicit narrowing to the
intersection of multi-objective optimisation and multi-energy
systems. The contribution is a structured synthesis along three
complementary axes -- a method taxonomy
(Section~\ref{sec:taxonomy}), an integration-pattern classification
(Section~\ref{sec:integration}), and a decision-aware validation
protocol (Section~\ref{sec:validation}) -- that together reframe how
surrogate-assisted MES studies should be designed and reported. The
application synthesis in Section~\ref{sec:applications} traces the
five integration patterns through dispatch and unit commitment,
optimal power flow, capacity expansion, district heating, multi-
energy systems, microgrids, multi-objective design, and stochastic
or robust planning, and shows that the dominant pattern shifts with
the underlying decision: P1 dominates inner-loop replacement
problems, P2 dominates expensive multi-objective design, P3 and P4
emerge in real-time and decomposition contexts, and P5 covers
chance-constrained and probabilistic planning.

Three observations cut across the screened literature. First, the
reported speed-ups are large and stable: one to three orders of
magnitude relative to truth-model evaluation are routinely achieved
across patterns and applications, and the principal source of gain is
the structural mismatch between the cost of an inner truth-model
evaluation and the number of times the host optimiser has to call it.
Second, the reported \emph{decision-quality} loss is far less
documented than the speed-up: regression metrics dominate the
reporting practice, decision-aware metrics (regret, MIP-gap
distribution, chance-constraint feasibility, calibrated tail
behaviour) appear only in a minority of studies. Third, surrogate-
assisted multi-objective design of multi-energy systems is the most
active subdomain of the recent literature; its open question has
shifted from \emph{whether} surrogates accelerate the Pareto search
to \emph{how} the resulting front, with its surrogate-induced
uncertainty, should be communicated to a decision maker.

The six research directions in Section~\ref{sec:challenges} -- a
decision-aware benchmark, calibrated uncertainty for surrogate-
assisted MOO, out-of-distribution generalisation across capacity
mixes, standardised public datasets, MCDM under surrogate
uncertainty and a shared open-source toolchain -- collectively define
a research agenda that aligns the methodological and the application
side. The methodological side has produced a rich vocabulary of
surrogate families and integration patterns; the application side
needs benchmarks, datasets and reporting conventions that allow
this vocabulary to be exercised under fair and reproducible
conditions. We hope that the taxonomy, the integration-pattern
classification, the decision-aware validation protocol and the
publicly released curated bibliography in this review provide a
concrete starting point for both sides.
